% =========================
% chapters/07_results_validation.tex
% =========================
\chapter{Results and Validation Templates}
\label{ch:results_validation}

\section{What ``validation'' means in this thesis}
Because no single metric captures ``good control,'' this thesis uses a small set of validation results:
(i) timing accuracy and repeatability, (ii) phase stability and repeatability, and (iii) spectral purity for single-tone
and multitone cases.

Each subsection below defines:
\emph{metric} $\rightarrow$ \emph{measurement method} $\rightarrow$ \emph{expected failure modes} $\rightarrow$ \emph{plots to include}.

\section{Timing accuracy}
\subsection{Metric}
Define the timing metric as the deviation between scheduled RTIO event time and observed analog event time.
Report both mean error and jitter (RMS).

\subsection{Method}
Use a scope with a common trigger derived from an ARTIQ TTL line scheduled at the same timestamp as the
Phaser update. Measure the relative delay and run-to-run variation over $N$ repetitions.

% TODO: Specify your instrument model and sampling rate.
\begin{figure}[t]
  \centering
  \draftfigure[0.95\linewidth]{figures/result_timing_histogram.pdf}
  \caption{Placeholder: histogram of timing error (scheduled vs observed) across trials.}
  \label{fig:result_timing_hist}
\end{figure}

\section{Phase stability}
\subsection{Metric}
For two channels (or two tones), define $\Delta\phi(t)$ as the relative phase. Report drift over time, and RMS
phase noise over defined bandwidths.

\subsection{Method options}
Option A: Use a vector signal analyzer (VSA) to extract instantaneous phase.
Option B: Digitize the waveform on a scope, perform IQ demodulation in software, and compute phase.
Option C: Heterodyne two channels into a beat note and analyze phase fluctuations of the beat.

\begin{figure}[t]
  \centering
  \draftfigure[0.95\linewidth]{figures/result_phase_drift.pdf}
  \caption{Placeholder: relative phase drift $\Delta\phi(t)$ over time under a fixed configuration.}
  \label{fig:result_phase_drift}
\end{figure}

\section{Spectrum purity and multitone behavior}
\subsection{Metrics}
Report SFDR for a single tone, and for multi-tone: intermodulation products and noise floor.
Document how performance depends on output power and tone count.

\subsection{Method}
Use a spectrum analyzer with sufficient dynamic range; ensure consistent RBW/VBW settings and correct
windowing/averaging. For multi-tone tests, document tone placement, spacing, and amplitude plan.

\begin{figure}[t]
  \centering
  \draftfigure[0.95\linewidth]{figures/result_spectrum_single_tone.pdf}
  \caption{Placeholder: spectrum of single-tone output with labeled harmonics and spurs.}
  \label{fig:result_spectrum_single}
\end{figure}

\begin{figure}[t]
  \centering
  \draftfigure[0.95\linewidth]{figures/result_spectrum_multitone.pdf}
  \caption{Placeholder: spectrum of multi-tone output with labeled intermodulation products.}
  \label{fig:result_spectrum_multi}
\end{figure}

\section{Example plot generation pipeline (template)}
% TODO: Either generate plots externally (Python/Matlab) and include PDFs,
% or use pgfplots with CSV input:
\begin{figure}[t]
  \centering
  \begin{tikzpicture}
    \begin{axis}[width=0.95\linewidth, xlabel={Frequency (Hz)}, ylabel={Power (dBc)}]
      %\addplot table [x=freq_hz, y=power_dbc, col sep=comma] {data/spectrum.csv};
    \end{axis}
  \end{tikzpicture}
  \caption{Placeholder: pgfplots-based spectrum figure (optional).}
  \label{fig:pgfplots_example}
\end{figure}
